\documentclass[a4paper,12pt]{article}
\usepackage[utf8]{inputenc}
\usepackage[T2A]{fontenc}
\usepackage[english, russian]{babel}
\usepackage{amsmath}
\usepackage{amssymb}
\usepackage{geometry}
\geometry{margin=2.5cm}
\usepackage{setspace}
\usepackage{titlesec}
\usepackage{indentfirst}
\usepackage{graphicx}
% Remove fontspec package
\usepackage{hyperref}
\hypersetup{
    colorlinks=true,
    linkcolor=blue,
    filecolor=magenta,
    urlcolor=cyan,
}

\begin{document}
\begin{center}
\small Министерство науки и высшего образования Российской Федерации\\
ФЕДЕРАЛЬНОЕ ГОСУДАРСТВЕННОЕ АВТОНОМНОЕ\\
ОБРАЗОВАТЕЛЬНОЕ УЧРЕЖДЕНИЕ ВЫСШЕГО ОБРАЗОВАНИЯ\\
НАЦИОНАЛЬНЫЙ ИССЛЕДОВАТЕЛЬСКИЙ УНИВЕРСИТЕТ\\
ИТМО
\vfill

\Large\textbf{Лабораторная работа №1}\\
\small{по дисциплине}\\
\small{«Численные методы анализа»}
\vfill
\thispagestyle{empty}

\begin{center}
Выполнили студенты:\\
\vspace{0.3cm}
Бессонов Александр Русланович\\
Группы J3113, ИСУ 465220\\
\vspace{0.3cm}
Тиганов Вадим Игоревич\\
Группы J3112, ИСУ 467701
\end{center}
\vfill

Санкт-Петербург\\
\the\year
\end{center}

\newpage



\tableofcontents


\newpage

\section{Цель работы:}
Изучение и применение методов сингулярного разложения (\verb|SVD|) для анализа
и восстановления изображений. 

\section{Постановка задачи:}
На основе выбранного изображения требуется провести следующие эксперименты:
\begin{enumerate}
    \item Преобразовать изображение в тензоры цветовых каналов \verb|R, G, B|
    \item Применение степенного метода для выполнения сингулярного разложения каждого цветового канала. Анализ точности восстановления изображения при различном числе итераций \verb|(10, 50, 100)| и сравнение результатов.
    \item Исследование распределения сингулярных чисел для каждого цветового канала и определение минимального числа старших сингулярных чисел, необходимых для приемлемого восстановления изображения.
    \item Нахождение сингулярных чисел одного из цветовых каналов с использованием метода вращений Якоби и сравнение результатов с результатами, полученными степенным методом. Оценка числа операций, выполненных в каждом методе.
\end{enumerate}

\newpage

\section{Описание данных:}
Было выбрано следующее изображение:

\begin{figure}[h]
    \centering
    \includegraphics[width=0.5\textwidth]{/home/vadim/Edu/Numerical-methods-labs/lab-2/data/pictures/photo.jpeg}
    \caption{Исходное изображение}
    \label{fig:original_image}
\end{figure}
Размер изображения: \verb|640x640| пикселей. Изображение состоит из трех цветовых каналов: \verb|R, G, B|.
\vspace{1cm}

Дополнительно визуализировали каждый цветовой канал:

\begin{figure}[h]
    \centering
    \includegraphics[width=1\textwidth]{/home/vadim/Edu/Numerical-methods-labs/lab-2/data/pictures/RGB-picture.png}
    \caption{Представление по каналам исходного изображения}
    \label{fig:original_image}
\end{figure}

Размеры изображения и его цветовых каналов при представлении в виде числового массива:
\begin{verbatim}
    Размер изображения: (640, 640, 3)
    Формы каналов: (640, 640), (640, 640), (640, 640)
\end{verbatim}


\section{Ход работы:}

\subsection{Преобразование изображения в тензоры цветовых каналов}
Использована библиотека \verb|opencv-python|
\begin{verbatim}
    image_bgr = cv.imread(path_to_image)
    R = image_bgr[:,:,0].astype(float)
    G = image_bgr[:,:,1].astype(float)
    B = image_bgr[:,:,2].astype(float)
\end{verbatim}
Числовое представление изображения представляет собой тензор размерности
\verb|(640,640,3)|, где каждый элемент соответствует тройке \verb|(R,G,B)| для каждого
пикселя слева направо и сверху вниз.\\
Каждый из тензоров цвета представляет собой матрицу размеров \verb|(640,640)|, где каждый элемент представляет собой
яркость каждого из цветовых каналов для каждого пикселя от \verb|0| до \verb|255|. 

\subsection{Сингулярное разложение с помощью степенного метода}
В этом разделе с помощью степенного метода (\texttt{power\_method\_svd}) приближённо вычисляется наибольшее сингулярное число $\sigma_1$ каждого цветового канала.

\subsubsection*{Алгоритм степенного метода}
Степенной метод позволяет найти наибольшее сингулярное число $\sigma_1$ и соответствующие левый и правый сингулярные векторы $u_1$ и $v_1$ для матрицы $A$.

Алгоритм основан на том, что если матрицу $A$ многократно применять к некоторому начальному вектору $v$, то результат будет стремиться к направлению доминирующего собственного вектора матрицы $A^TA$.

\textbf{Шаги алгоритма:}
\begin{enumerate}
    \item Инициализируем случайный единичный вектор $v \in \mathbb{R}^n$
    \item На каждой итерации вычисляем:
    \begin{align}
        w &= A^T(Av) \\
        v_{new} &= \frac{w}{\|w\|}
    \end{align}
    \item Проверяем условие сходимости: $\|v_{new} - v\| < \varepsilon$, где $\varepsilon$ — заданная точность
    \item После сходимости вычисляем:
    \begin{align}
        u &= Av \\
        \sigma &= \|u\| \\
        u &= \frac{u}{\sigma}
    \end{align}
\end{enumerate}

В результате получаем наибольшее сингулярное число $\sigma$ и соответствующие ему левый и правый сингулярные векторы $u$ и $v$.

% Канал R:
% N=10: 	Ошибка = 32975.45
% N=50: 	Ошибка = 32975.45
% N=100: 	Ошибка = 32975.45
% Канал G:
% N=10: 	Ошибка = 19765.48
% N=50: 	Ошибка = 19765.48
% N=100: 	Ошибка = 19765.48
% Канал B:
% N=10: 	Ошибка = 54508.41
% N=50: 	Ошибка = 54508.41
% N=100: 	Ошибка = 54508.41
\subsubsection*{Результаты применения степенного метода}

% Для оценки точности приближения строится ранг-1 матрица:
% $$
% \tilde{A} = \sigma_1 \cdot u_1 v_1^T
% $$
% и вычисляется евклидова ошибка восстановления:
% $$
% \text{Ошибка} = \left\| A - \tilde{A} \right\|_2
% $$
Для оценки точности приближения строится ранг-1 матрица:
$$
\tilde{A} = \sigma_1 \cdot u_1 v_1^T
$$
и вычисляется ошибка восстановления с помощью нормы Фробениуса:
$$
\text{Ошибка} = \left\| A - \tilde{A} \right\|_F = \sqrt{\sum_{i=1}^{m}\sum_{j=1}^{n} \left(A_{ij} - \tilde{A}_{ij}\right)^2}
$$
где $A$ - исходный цветовой канал изображения, а $\tilde{A}$ - его приближение ранга 1, полученное с помощью степенного метода. Норма Фробениуса представляет собой корень из суммы квадратов всех элементов матрицы разности и является естественной мерой отклонения между матрицами.
Для оценки точности и сходимости алгоритма мы измеряли ошибку восстановления при различном числе итераций. Результаты представлены в таблице \ref{tab:power_method_error}.

\begin{table}[h]
\centering
\begin{tabular}{|c|c|c|c|}
\hline
\textbf{Цветовой канал} & \textbf{10 итераций} & \textbf{50 итераций} & \textbf{100 итераций} \\
\hline
R & 32975.45 & 32975.45 & 32975.45 \\
\hline
G & 19765.48 & 19765.48 & 19765.48 \\
\hline
B & 54508.41 & 54508.41 & 54508.41 \\
\hline
\end{tabular}
\caption{Ошибка восстановления изображения при различном числе итераций}
\label{tab:power_method_error}
\end{table}

Как видно из таблицы, ошибка восстановления не изменяется при увеличении числа итераций с 10 до 100. Это может говорить о том, что метод быстро сходится на данном изображении, и уже после 10 итераций дальнейшего улучшения не происходит.

\subsection{Распределение сингулярных чисел и усечённое SVD}
Использовали метод \verb|numpy.linalg.svd| для вычисления сингулярных чисел каждого цветового канала. Сингулярные числа отсортированы по убыванию.\\
Для каждого цветового канала вычисляем сингулярные числа и строим график их распределения.\\
\begin{verbatim}
_, sR, _ = np.linalg.svd(R, full_matrices=False)
_, sG, _ = np.linalg.svd(G, full_matrices=False)
_, sB, _ = np.linalg.svd(B, full_matrices=False)
print(len(sR), len(sG), len(sB))  
\end{verbatim}


\begin{figure}[h]
    \centering
    \includegraphics[width=0.7\textwidth]{/home/vadim/Edu/Numerical-methods-labs/lab-2/data/pictures/singular-values.png}
    \caption{Распределение сингулярных чисел цветовых каналов}
    \label{fig:singular_values}
\end{figure}

\subsubsection{Анализ распределения сингулярных чисел}

Анализируя графики распределения сингулярных чисел (рис. \ref{fig:singular_values}), можно сделать следующие наблюдения:

\begin{itemize}
    \item \textbf{Стремительный спад в начале}: Первые $50$–$100$ сингулярных чисел содержат подавляющую часть энергии изображения. На графике наблюдается резкое падение от $i=1$ до $i\approx100$.
    
    \item \textbf{Сходство форм кривых}: Для всех трех цветовых каналов (R, G, B) формы распределений практически идентичны, что говорит о схожей структуре информации во всех каналах. При этом канал B демонстрирует несколько более высокие значения $\sigma_i$ в начале распределения, что может отражать более высокую контрастность синих компонент в исходном изображении.
    
    \item \textbf{Длинный «хвост»}: После $i \approx 300$ сингулярные числа падают ниже значения $0.1$, что указывает на то, что оставшиеся компоненты содержат преимущественно шум или крайне мелкие детали изображения.
\end{itemize}

Анализ спектра показывает, что оптимальное соотношение между объёмом данных и качеством восстановления изображения достигается при $k$ в диапазоне $50$–$100$. Для практического сжатия достаточно усечённого SVD с $k \approx 50$, а для почти безупречного восстановления целесообразно увеличить $k$ до $100$–$150$.

\subsubsection{Восстановление изображения с помощью усечённого SVD}
Для восстановления изображения с помощью усечённого SVD необходимо оставить только первые $k$ сингулярных чисел и соответствующие им сингулярные векторы. Восстановление производится по формуле:
$$A_k = \sum_{i=1}^{k} \sigma_i u_i v_i^T = U_k \Sigma_k V_k^T$$

где:
\begin{itemize}
    \item $A_k$ — аппроксимация исходной матрицы $A$ с рангом $k$
    \item $U_k$ — матрица из первых $k$ левых сингулярных векторов
    \item $\Sigma_k$ — диагональная матрица из первых $k$ сингулярных чисел
    \item $V_k^T$ — транспонированная матрица из первых $k$ правых сингулярных векторов
\end{itemize}

Проведено восстановление изображения с различным числом сингулярных компонент $k$. На рисунке \ref{fig:reconstructed_images} представлены результаты реконструкции для $k = 10$, $k = 50$ и $k = 100$.

\begin{figure}[h]
    \centering
    \includegraphics[width=1\textwidth]{/home/vadim/Edu/Numerical-methods-labs/lab-2/data/pictures/reconstructed-images.png}
    \caption{Восстановленное изображение при различных значениях $k$}
    \label{fig:reconstructed_images}
\end{figure}

Визуальный анализ показывает:
\begin{itemize}
    \item При $k = 10$ отчетливо видны только основные контуры и цветовые области, мелкие детали отсутствуют
    \item При $k = 50$ большинство значимых деталей уже восстановлено, но текстуры все еще сглажены
    \item При $k = 100$ изображение визуально практически неотличимо от оригинала
\end{itemize}

Степень сжатия данных при усечённом SVD можно оценить как отношение объёма хранимых данных в оригинальном и сжатом изображении:
$$\text{Коэффициент сжатия} = \frac{m \times n}{k \times (m + n + 1)}$$

где $m \times n$ — размер исходной матрицы, а $k \times (m + n + 1)$ — количество элементов, необходимых для хранения $k$ сингулярных чисел и соответствующих векторов.


\subsection{Вычисление сингулярных чисел методом вращений Якоби}

\subsubsection{Теоретическое обоснование}

Метод вращений Якоби — это итеративный алгоритм для нахождения собственных значений и собственных векторов симметричной матрицы. В контексте SVD мы применяем его к матрице $A^TA$, чтобы найти собственные значения, квадратные корни которых дают нам сингулярные числа матрицы $A$.

\textbf{Принцип метода:} На каждой итерации выбирается максимальный внедиагональный элемент $a_{pq}$ и применяется вращение в плоскости $(p,q)$, обнуляющее этот элемент.

Матрица вращения $J(p,q,\theta)$ имеет вид:
\begin{align}
J_{ij} = 
\begin{cases}
\cos\theta, & \text{если } i=j=p \text{ или } i=j=q \\
\sin\theta, & \text{если } i=p, j=q \\
-\sin\theta, & \text{если } i=q, j=p \\
1, & \text{если } i=j \neq p,q \\
0, & \text{в остальных случаях}
\end{cases}
\end{align}

Угол вращения $\theta$ выбирается так, чтобы обнулить элемент $a_{pq}$:
\begin{align}
\tan(2\theta) = \frac{2a_{pq}}{a_{pp} - a_{qq}}
\end{align}

При $a_{pp} = a_{qq}$ угол $\theta = \pi/4$.

\subsubsection*{Процесс алгоритма:}
\begin{enumerate}
    \item Находим максимальный по модулю внедиагональный элемент $a_{pq}$
    \item Вычисляем угол вращения $\theta$
    \item Применяем матрицу вращения: $A' = J^T A J$
    \item Повторяем процесс до достижения заданной точности (все внедиагональные элементы меньше $\varepsilon$)
\end{enumerate}

После сходимости диагональные элементы матрицы являются собственными значениями, а столбцы накопленной матрицы вращений — собственными векторами.

\subsubsection{Сравнение методов}

Для сравнения методов мы вычислили сингулярные числа канала R тремя способами:
\begin{itemize}
    \item Методом вращений Якоби
    \item Степенным методом с дефляцией
    \item \verb|numpy| методом \verb|linalg.svd|
\end{itemize}

\begin{figure}[h]
    \centering
    \includegraphics[width=0.8\textwidth]{/home/vadim/Edu/Numerical-methods-labs/lab-2/data/pictures/singular-values-compare.png}
    \caption{Сравнение сингулярных чисел, полученных методом Якоби и степенным методом}
    \label{fig:singular_values_comparison}
\end{figure}
\subsubsection{Анализ результатов сравнения методов}

Из графика (рис. \ref{fig:singular_values_comparison}) можно сделать следующие выводы:

\begin{itemize}
    \item \textbf{Точность вычислений:} 
    Наблюдается существенное расхождение между методами, со средней абсолютной ошибкой порядка $2.5 \times 10^{3}$ и нормой разности векторов сингулярных чисел $4.8 \times 10^{4}$.
    
    \item \textbf{Особенности метода Якоби:} 
    Метод корректно определяет только первые несколько значений $\sigma_i$, после чего кривая практически выравнивается и не отражает экспоненциальный спад, характерный для истинного спектра сингулярных чисел.
    
    \item \textbf{Поведение степенного метода:} 
    Степенной метод с дефляцией более точно воспроизводит плавно убывающий профиль, характерный для сингулярных чисел изображения.
    
    \item \textbf{Эталонное решение:} 
    Реализация \texttt{numpy.linalg.svd} предоставляет наиболее точный результат, который мы принимаем за эталонный при оценке других методов.
\end{itemize}

\subsection{Гистограммы глубины цвета и статистика}

Был проведён анализ распределения интенсивностей для каждого цветового канала. На рисунке \ref{fig:color_histograms} представлены гистограммы и основные статистические показатели.

\begin{figure}[h]
    \centering
    \includegraphics[width=1.1\textwidth]{/home/vadim/Edu/Numerical-methods-labs/lab-2/data/pictures/chan-analysis.png}
    \caption{Гистограммы и статистические показатели цветовых каналов}
    \label{fig:color_histograms}
\end{figure}

\subsubsection*{Анализ распределения интенсивностей}

Из представленных гистограмм и статистических данных можно сделать следующие выводы:

\begin{itemize}
    \item \textbf{Высокая дисперсия во всех каналах:} Стандартное отклонение превышает 68 для всех трёх каналов, что свидетельствует о высоком контрасте и визуальной насыщенности изображения.
    
    \item \textbf{Особенности R-канала:} Красный канал имеет наиболее выраженные крайние значения как в тенях, так и в светлых участках, что делает его наиболее насыщенным по интенсивности.
    
    \item \textbf{Характеристики B-канала:} Синий канал наименее выражен, его более низкие значения указывают на преобладание тёплых тонов, что характерно для фотографий с естественным освещением или снимков с тёплым цветовым балансом.
\end{itemize}

Широкие распределения и множественные пики на гистограммах указывают на высокую информационную плотность изображения, что подтверждает необходимость использования достаточно большого числа сингулярных чисел для качественной реконструкции.

\newpage

\section{Выводы и заключение:}
В ходе данной лабораторной работы были получены следующие результаты:

\begin{enumerate}
    \item \textbf{Загрузка и разложение каналов изображения}
    \begin{itemize}
        \item Цветное изображение разбивалось на три матрицы каналов $R$, $G$, $B$.
        \item Для каждого из них реализованы и применены три подхода:
        \begin{itemize}
            \item Степенной метод с дефляцией (усечение SVD).
            \item Полный SVD через NumPy.
            \item Метод вращений Якоби для спектра матрицы $A^T A$.
        \end{itemize}
    \end{itemize}
    
    \item \textbf{Степенной метод (ранг-1)}
    \begin{itemize}
        \item Для ранга-1 приближения оценили погрешность $\|A - \sigma_1 u_1 v_1^T\|_F$.
        \item Ошибка стабилизировалась уже при $N \approx 10$ итерациях для всех каналов, что указывает на быструю сходимость первой компоненты.
    \end{itemize}
    
    \item \textbf{Полный спектр сингулярных чисел}
    \begin{itemize}
        \item Графики $\{\sigma_i\}$ показали, что первые $50$--$100$ чисел содержат основную «энергию» изображения.
        \item Последующий «хвост» при $i > 300$ практически стремится к нулю и отвечает преимущественно за шум и мельчайшие детали.
    \end{itemize}
    
    \item \textbf{Усечённый SVD (ранг-$k$)}
    \begin{itemize}
        \item При $k=10$ восстановление даёт сильное размытие и сохраняет лишь основные цветовые области.
        \item При $k=50$ достигается оптимальный компромисс между качеством и объёмом хранимых данных.
        \item При $k=100$ визуальное отличие от оригинала становится минимальным, однако объём хранения растёт пропорционально $k$.
    \end{itemize}
    
    \item \textbf{Сравнение методов Якоби и степенного}
    \begin{itemize}
        \item Алгоритм вращений Якоби с параметрами $\text{max\_iter}=100$, $\text{tol}=10^{-10}$ даёт точные первые несколько $\sigma_i$, после чего наблюдается существенное отклонение от эталонного решения.
        \item Степенной метод с дефляцией лучше воспроизводит плавно убывающий профиль сингулярных чисел, согласующийся с результатами NumPy-SVD.
    \end{itemize}
\end{enumerate}

В целом, сингулярное разложение показало себя эффективным методом для анализа и сжатия изображений. Удалось определить, что для качественного восстановления данного изображения достаточно использовать около 100 компонент из 640 возможных, что даёт значительную степень сжатия при сохранении визуального качества.



\end{document}
